\usepackage{caption}
\usepackage{amsmath,amsfonts, amssymb}
\usepackage{subfigure}
\usepackage{hyperref}
\usepackage{url}
\def\UrlBreaks{\do\A\do\B\do\C\do\D\do\E\do\F\do\G\do\H\do\I\do\J\do\K\do\L\do\M\do\N\do\O\do\P\do\Q\do\R\do\S\do\T\do\U\do\V\do\W\do\X\do\Y\do\Z\do\[\do\\\do\]\do\^\do\_\do\`\do\a\do\b\do\c\do\d\do\e\do\f\do\g\do\h\do\i\do\j\do\k\do\l\do\m\do\n\do\o\do\p\do\q\do\r\do\s\do\t\do\u\do\v\do\w\do\x\do\y\do\z\do\.\do\@\do\\\do\/\do\!\do\_\do\ \do\;\do\>\do\]\do\)\do\,\do\?\do\'\do+\do\=\do\#}
\usepackage{graphicx}
\usepackage{textcomp}
\usepackage{xcolor}
\usepackage{amsthm}
\usepackage{colortbl}
\usepackage[misc]{ifsym}
\usepackage{xspace}
\usepackage{multirow}
\usepackage{threeparttable}
\usepackage{subfigure}
\usepackage{booktabs}
\usepackage{tabularx}
\let\labelindent\relax
\usepackage{enumitem}
\usepackage{booktabs}
\usepackage[linesnumbered,ruled,vlined]{algorithm2e}
\usepackage[normalem]{ulem}
\usepackage{dblfloatfix}	
\usepackage[euler]{textgreek}

% \usepackage{mdframed}
\usepackage{framed}
\definecolor{formalshade}{rgb}{0.95,0.95,0.97}
\definecolor{darkblue}{rgb}{0.14,0.22,0.52}
\newenvironment{formal}{
  \def\FrameCommand{
    \hspace{0pt}
    {\color{darkblue}\vrule width 1.5pt}
    \colorbox{formalshade}
  }
  \MakeFramed{\advance\hsize-\width\FrameRestore}
  \noindent\hspace{-4.55pt}
  % \begin{adjustwidth}{}{7pt}
}
{
  % \end{adjustwidth}
  \endMakeFramed%
}

\SetKwInput{KwInput}{Input}                % Set the Input
\SetKwInput{KwOutput}{Output}              % set the Output

\newtheorem{theorem}{Theorem}
\newtheorem{lemma}[theorem]{Lemma}
\newtheorem{example}{Example}
\newtheorem{definition}{Definition}
\newtheorem{remark}{Remark}
\newtheorem{property}{Property}
\DeclareMathOperator{\E}{\mathbb{E}}

\def\eg{\emph{e.g.,}\xspace}
\def\etc{\emph{etc.,}\xspace}
\def\ie{\emph{i.e.,}\xspace}
\def\etal{\emph{et al.}\xspace}
\def\vs{\emph{vs.}\xspace}
\def\cf{\emph{cf.}\xspace}

\newcommand{\one}{({\em i}\/)\xspace}
\newcommand{\two}{({\em ii}\/)\xspace}
\newcommand{\three}{({\em iii}\/)\xspace}
\newcommand{\four}{({\em iv}\/)\xspace}
\newcommand{\five}{({\em v}\/)\xspace}
\newcommand{\six}{({\em vi}\/)\xspace}
\newcommand{\seven}{({\em vii}\/)\xspace}

\newcommand{\todo}[1]{}
\renewcommand{\todo}[1]{{\color{red} TODO: {#1}}}
\newcommand\mycommfont[1]{\footnotesize\ttfamily\textcolor{blue}{#1}}
\SetCommentSty{mycommfont}

\usepackage{cleveref}
\crefformat{section}{\S\textcolor{purple}{#2#1#3}} 
\crefformat{subsection}{\S\textcolor{purple}{#2#1#3}}
\crefformat{subsubsection}{\S\textcolor{purple}{#2#1#3}}
\crefformat{figure}{#2{\color{blue}#1}#3}
\crefformat{table}{#2{\color{blue}#1}#3}
\crefformat{appendix}{Appendix \textcolor{purple}{#2#1#3}}

\newcommand{\tool}{\textsc{CAUSE}\xspace}

\pagenumbering{gobble}

\hyphenation{op-tical net-works semi-conduc-tor}

\newcommand{\ruoxi}[1]{\textcolor{violet}{Ruoxi: #1}}

\newcommand{\jason}[1]{\textcolor{red}{Jason: #1}}

\newcommand{\update}[1]{\textcolor{blue}{#1}}